\section{Introducción}

El presente proyecto de investigación tiene como finalidad principal valorar si la estructuración del desarrollo de sintetizadores de sonido mediante la aplicación de \textit{Model-Driven Architecture} (MDA), adaptando el modelado a las necesidades del dominio de la música electrónica, y orientándolo a la generación automática de código a través de la integración con herramientas de IA, puede aportar mejoras significativas en el desarrollo de este tipo de software. Para lograrlo, las herramientas desarrolladas se centrarán en las dos primeras capas de abstracción que propone MDA: el \textit{Computation Independent Model} (CIM) y el \textit{Platform Independent Model} (PIM).

Con el objetivo de orientar y formalizar estos desarrollos, se han diseñado e implementado una serie de lenguajes de dominio específicos (DSL) al ámbito en cuestión, 2 para el modelado CIM y otros 2 para el modelado PIM. Estos lenguajes permiten capturar la semántica y la arquitectura de los desarrollos en 2 diferentes niveles de abstracción. Como novedad, respecto al uso convencional de la capa CIM por otras herramientas, donde la escritura de los requisitos se representa de manera desestructurada, en este trabajo se propone una estructuración cerrada y estricta para esta capa, pero sin perder el nivel de abstracción que requiere esta capa. 

Una de las grandes ventajas de esta estructuración basada en modelos estrictos para las capas CIM / PIM, y su relacionabilidad entre ellas, es tanto la capacidad de generar \textit{prompts} altamente estructurados y fáciles de procesar por inteligencias artificiales como la de poder aplicar técnicas de ingeniería inversa entre los distintos modelos desarrollados. Esta integración brinda la posibilidad de iterar de manera continua y eficiente entre las fases de diseño y análisis de los modelos desarrollados, permitiendo obtener con ello, de forma iterativa, nuevas posibles generaciones del código fuente objetivo.

Para materializar este enfoque, se ha desarrollado una herramienta integral, llamada MDAA, que implementa, a modo de IDE de desarrollo, todo el proceso de análisis y diseño para el dominio ejemplo, cuyo código fuente se encuentra en la carpeta \textit{\_CODIGO/} del paquete entregado, y también en el repositorio GitHub\footnote{\url{\urlProyecto}} del proyecto. Además, aprovechando el acceso gratuito para estudiantes de UNED en Azure, se puede acceder a ella de forma remota a través de enlace web\footnote{\url{\urlProyectoMdaa}} (horario de activación, de 8:00 a 20:00, de lunes a viernes), usando el usuario y contraseña \texttt{test@uned.es} / \texttt{11111111}. Esta plataforma ofrece un entorno completo para la definición visual de los modelos, facilitando el ciclo de vida del desarrollo y comprobación de las tesis sostenidas por este artículo. Los principales puntos que ofrece la herramienta son:

\begin{itemize}
    \item \textbf{Gestión de modelos CIM y PIM:} Entornos dedicados para la implementación de los modelos de análisis y diseño. Desarrollo gráfico de los modelos PIM, facilitando su implementación mediante la abstracción de los aspectos de bajo nivel postulados por los DSL desarrollados.
    \item \textbf{Lenguajes DSL especializados:} Integración de gramáticas propias que validan la coherencia estructural de las conexiones y flujos de señales que debe tener una aplicación software de estas características.
    \item \textbf{Exportación asistida para IA:} Capacidad de exportar el contexto del proyecto junto con los manuales de los DSL, adaptándolos de cara a generar un prompt estructurado y coherente, orientado a guiar a una IA en la generación, de forma estricta, del código deseado.
    \item \textbf{Validación continua:} Verificación en tiempo real para garantizar que los modelos cumplen con las restricciones del dominio antes de su implementación.
\end{itemize}

Aparte, en el mismo repositorio de la herramienta, en la carpeta \textit{\_CODIGO/dsl}, se puede encontrar información detallada de los modelos DSL desarrollados y su código. Los DSL desarrollados constituyen el núcleo de este proyecto de investigación, intentando buscar con ellos una forma de estructurar desarrollos en este dominio concreto.

En el repositorio entregado, disponemos tanto de un manual completo de la aplicación IDE mencionada, MDAA\footnote{\url{\urlProyectoManualMdaa}}, como de los DSL\footnote{\url{\urlProyectoManualDsl}} desarrollados.

\textbf{Nota importante:} Todo documento, enlace de código o información mencionados en este artículo como accesible bajo el repositorio del proyecto\footnote{\url{\urlProyecto}}, se encuentra disponible también, bajo la misma ruta, en el .zip adjunto.

\textbf{Recomendación:} Debido a la especificidad de la materia, se recomienda consultar el manual complementario\footnote{\url{\urlProyectoMaterialComplementario}} sobre síntesis de sonido y Pure Data. Este documento, realizado ex profeso para lectores no familiarizados con el dominio abordado en este artículo, ofrece una visión general sobre la síntesis de sonido digital y unas nociones básicas de lo que es \textit{Pure Data}, el lenguaje utilizado en las pruebas.

